\documentclass[12pt]{article}
\usepackage[utf8]{inputenc}
\usepackage[T1]{fontenc}
\usepackage[margin=1in]{geometry}
\usepackage{setspace}
\usepackage{xcolor}
\usepackage[colorlinks=true, urlcolor=blue!60!black, linkcolor=blue!60!black]{hyperref}
\setlength{\parindent}{2em}
\setlength{\parskip}{1.25em}
\renewcommand{\baselinestretch}{1.0}
\begin{document}
\begin{center}
{\LARGE \textit{Church of the Lukumi Babalu Aye, Inc. v. City of Hialeah}, 508 U.S. 520 (1993)\par}
\vspace{0.8em}
{\large Case Brief --- Four-Page Edition\par}
\end{center}
\vspace{1em}

\section{Memory Jogger}

Municipal ordinances gerrymandered to ban Santeria ritual animal sacrifice while exempting almost all secular and kosher animal killings are neither neutral nor generally applicable under \textit{Employment Division v. Smith}, and therefore trigger---and fail---strict scrutiny under the Free Exercise Clause (\href{https://supreme.justia.com/cases/federal/us/508/520/}{508 U.S. 520}).

\section{Detailed Case Facts}

The Church of the Lukumi Babalu Aye and its priest, Ernesto Pichardo, practiced Santeria, an Afro-Caribbean religion fusing Yoruba worship of \textit{orishas} with Catholic iconography, brought to South Florida by Cuban exiles. Animal sacrifice---killing chickens, goats, and other animals by cutting the carotid arteries, usually followed by eating the meat---is central to its rites. In 1987 the Church leased land in Hialeah and announced plans to open a house of worship, school, and cultural center, bringing its rituals into the open.

At an emergency council session marked by open hostility, council members said Santeria adherents ``are in violation of everything this country stands for,'' and the police chaplain urged the council ``not to permit this Church to exist.'' Within three months the council unanimously passed two resolutions and four ordinances. They incorporated Florida's ban on ``unnecessary'' killing (which the state Attorney General read to reach religious sacrifice), defined ``sacrifice'' as killing an animal in a ritual ``not for the primary purpose of food consumption,'' banned sacrifice within city limits, and confined slaughter to slaughterhouse zones while exempting licensed food establishments (thus kosher slaughter) and small numbers of hogs and cattle.

\section{Procedural History}

The Church sued under 42 U.S.C. \S~1983. After a nine-day bench trial, the District Court acknowledged that the ordinances were not religiously neutral but upheld them as serving compelling interests in public health, child welfare, animal welfare, and zoning. The Eleventh Circuit affirmed in a one-paragraph \textit{per curiam} opinion without evaluating \textit{Smith}. The Supreme Court granted certiorari and unanimously reversed on June 11, 1993.

\section{Judicial Votes}

Reversed, \textbf{9--0}, though the Justices divided over legislative motive and \textit{Smith}.
\begin{itemize}
\item \textbf{Opinion of the Court (Kennedy):} Parts I, III, and IV for seven Justices; Part II-B (general applicability) for six; Parts II-A-1 and II-A-3 (neutrality) for five; Part II-A-2 (legislative motive) joined only by Stevens.
\item \textbf{Concurring in part and in the judgment:} Scalia, joined by Rehnquist (all but Part II-A-2); Souter (Parts I, III, and IV only).
\item \textbf{Concurring in the judgment:} Blackmun, joined by O'Connor.
\end{itemize}

\section{Holding}

The four ordinances are unconstitutional. A law burdening religious practice need not satisfy strict scrutiny if it is neutral and generally applicable. But a law that is not neutral (its object is to restrict practices because of their religious motivation, overtly or through a ``religious gerrymander'') or not generally applicable (it pursues its interests only against religious conduct) must advance interests ``of the highest order'' and be narrowly tailored. The restatement of the \textit{Smith} safe harbor was arguably dicta, Part II-A-2's motive inquiry drew only two votes, and the Court declined to define the outer boundary of general applicability.

\section{Analysis of Opinions}

\textbf{Kennedy (Court).} Facial neutrality is only the minimum: the Clause forbids ``subtle departures from neutrality'' and ``religious gerrymanders.'' The interlocking definitions and exemptions made Santeria sacrifice virtually the only conduct banned, while hunting, fishing, pest extermination, and euthanasia remained lawful; the ordinances also swept far more broadly than public health required. They were substantially underinclusive as to both animal cruelty and public health, imposing on Santeria worshippers a burden society would not impose on itself. Under strict scrutiny they were neither narrowly tailored nor backed by compelling interests, since a law leaving ``appreciable damage'' to an interest unprohibited cannot protect an interest of the highest order. In Part II-A-2, Kennedy and Stevens relied on \textit{Arlington Heights} and the council's statements to find that the laws targeted Santeria.

\textbf{Scalia (with Rehnquist).} Neutrality and general applicability overlap. Courts should not probe lawmakers' subjective motives, which are ``virtually impossible'' to determine; the First Amendment addresses what laws do, not why they were passed.

\textbf{Souter.} Joining only the nonpersecution holding, Souter attacked \textit{Smith} as resting on weak distinctions, decided without full briefing, and at odds with the Clause's text and history, and urged reexamination.

\textbf{Blackmun (with O'Connor).} \textit{Smith} was wrong; any substantial burden on religion requires strict scrutiny, and a law that targets religion fails it \textit{ipso facto}. A genuinely neutral animal-cruelty law would be a harder case.

\section{Examples: Future Applications}

\begin{enumerate}
\item \textbf{Same Side (Unconstitutional): Targeting Halal Slaughter.} A city bans ``ceremonial animal killing'' in response to a mosque's Eid plans while exempting abattoirs, hunting, and backyard poultry slaughter. The gerrymandered exemptions fail neutrality and general applicability.
\item \textbf{Same Side (Unconstitutional): Worship Limits in an Emergency.} A governor caps in-home worship at three households while letting salons, stores, and restaurants host larger crowds. Favoring comparable secular activity triggers strict scrutiny, as in \textit{Tandon v. Newsom}.
\item \textbf{Opposite Side (Constitutional): Endangered-Bird Ban.} An across-the-board ban on possessing endangered birds' feathers, with no secular exemptions, is neutral and generally applicable and receives only rational-basis review.
\item \textbf{Opposite Side (Constitutional): Uniform Waste Rule.} A city requires everyone, from restaurants to houses of worship, to seal organic animal waste in leak-proof containers---the narrower regulation Kennedy said Hialeah could have adopted.
\item \textbf{Fence-Sitter: Hostile Commissioners.} An agency enforces a neutral public-accommodations law against a religious baker after commissioners disparage his faith, as in \textit{Masterpiece Cakeshop}. Under Part II-A-2, hostility taints neutrality; under Scalia's view, only the law's text and enforcement matter.
\end{enumerate}

\section{Critique}

\textbf{Doctrinal gaps.} Because Hialeah's targeting was so egregious, the Court never defined how much secular underinclusiveness defeats general applicability, leaving lower courts divided over the ``most-favored-nation'' theory until \textit{Fulton} and \textit{Tandon}. Part II-A-2's two votes left uncertain when officials' hostile remarks invalidate a law, and, as Blackmun noted, strict scrutiny after a finding of persecution is arguably redundant.

\textbf{Competing camps.} Textualists praise Scalia's refusal to psychoanalyze legislatures, while originalist critics of \textit{Smith}, led by Michael McConnell, argue that \textit{Lukumi} protects minorities only from clumsy or openly hostile majorities. Progressives celebrate the shield for minority faiths but warn that its general-applicability language has become a deregulatory sword against civil-rights and public-health laws.

\section{Key Quotations}

\begin{quote}
``Official action that targets religious conduct for distinctive treatment cannot be shielded by mere compliance with the requirement of facial neutrality. The Free Exercise Clause protects against governmental hostility which is masked as well as overt.''
\end{quote}
\noindent --- Justice Kennedy, Opinion of the Court, 508 U.S. at 534.

\begin{quote}
``The Free Exercise Clause commits government itself to religious tolerance, and upon even slight suspicion that proposals for state intervention stem from animosity to religion or distrust of its practices, all officials must pause to remember their own high duty to the Constitution and to the rights it secures.''
\end{quote}
\noindent --- Justice Kennedy, Opinion of the Court, 508 U.S. at 547.

\begin{quote}
``The compelling interest standard that we apply once a law fails to meet the \textit{Smith} requirements is not `water[ed] \dots{} down' but `really means what it says.'\,''
\end{quote}
\noindent --- Justice Kennedy, Opinion of the Court, 508 U.S. at 546.

\end{document}
